% 第四章：屈曲分析公式推导
% 包含：混合离散屈曲控制方程、能量方程、弱形式

本章针对 Mindlin 中厚板的屈曲问题，基于第二章建立的运动学方程，推导混合离散格式下的屈曲控制方程与广义特征值问题。首先采用 Green-Lagrange 应变张量推导屈曲问题的虚功方程；其次进行离散化，得到材料刚度矩阵与几何刚度矩阵的各子块表达式；最后讨论混合离散格式的稳定性条件。

\section{混合离散屈曲控制方程与特征值问题}

考虑 Mindlin 中厚板的屈曲问题。板在面内压力作用下，当荷载增大至临界值时，原始平面内平衡状态失去稳定性，板突然发生横向挠度（屈曲分岔）。与强度失效不同，屈曲失效通常发生在应力远低于材料屈服强度时，具有突发性。求解目标为临界荷载系数 $\lambda$ 和屈曲模态，对应的广义特征值问题为：
\begin{equation}
(\mathbf{K} + \lambda \mathbf{K}_G)\boldsymbol{\phi} = \mathbf{0}
\end{equation}
式中，$\mathbf{K}$ 为弹性刚度矩阵，$\mathbf{K}_G$ 为几何刚度矩阵，$\lambda$ 为屈曲荷载因子，$\boldsymbol{\phi}$ 为屈曲模态。临界屈曲荷载为 $P_{cr} = \lambda_{cr} P_{ref}$，其中 $P_{ref}$ 为参考预加荷载，$\lambda_{cr}$ 为最小正特征值。

面内位移变分 $u_\alpha$ 不会与预应力度 $\sigma_{\alpha\beta}$ 产生导致屈曲的几何力矩，因此几何刚度矩阵中与 $u_\alpha$ 对应的行列皆为零。实际控制板分岔屈曲的广义特征值问题子块矩阵方程为：
\begin{equation}
(\mathbf{k}_b + \lambda \mathbf{k}_g)\mathbf{b} = \mathbf{0}
\end{equation}

\section{屈曲分析能量方程}

由第二章建立的 Mindlin 板位移场可知，屈曲分析需考虑几何非线性，采用完整的 Green-Lagrange 应变张量：
\begin{equation}
\varepsilon_{ij} = \frac{1}{2}(\bar{u}_{i,j} + \bar{u}_{j,i} + \bar{u}_{k,i}\bar{u}_{k,j})
\end{equation}

将位移场代入，面内应变分量 $\varepsilon_{\alpha\beta}$ 按 $x_3$ 的幂次分项归类为：
\begin{equation}
\begin{aligned}
\varepsilon_{\alpha\beta} =& \underbrace{\frac{1}{2}(u_{\alpha,\beta} + u_{\beta,\alpha} + u_{\gamma,\alpha}u_{\gamma,\beta})}_{\text{薄膜应变}} + \underbrace{\frac{1}{2}w_{,\alpha}w_{,\beta}}_{\text{von K\'arm\'an 项}} \\
&- x_3\underbrace{\frac{1}{2}(\varphi_{\alpha,\beta} + \varphi_{\beta,\alpha} + u_{\gamma,\alpha}\varphi_{\gamma,\beta} + \varphi_{\gamma,\alpha}u_{\gamma,\beta})}_{\text{广义曲率}} \\
&+ x_3^2\underbrace{\frac{1}{2}\varphi_{\gamma,\alpha}\varphi_{\gamma,\beta}}_{\text{高阶旋转项}}
\end{aligned}
\end{equation}

横向剪切应变 $\varepsilon_{\alpha 3}$ 和厚度方向法向应变 $\varepsilon_{33}$ 分别为：
\begin{equation}
\varepsilon_{\alpha 3} = \frac{1}{2}(w_{,\alpha} - \varphi_\alpha - u_{\beta,\alpha}\varphi_\beta) + \frac{x_3}{2}\varphi_{\beta,\alpha}\varphi_\beta, \qquad \varepsilon_{33} = \frac{1}{2}\varphi_\alpha\varphi_\alpha
\end{equation}
上述各应变分量的完整逐项展开过程参见附录 B。

根据虚功原理，板的内虚功等于外虚功：
\begin{equation}
\int_{-h/2}^{h/2} \varepsilon_{ij}\sigma_{ij}\, dx_3\, d\Omega = \delta W_{ext}
\end{equation}
将应变张量按面内分量、横向剪切分量和厚度方向分量拆分为三部分积分：
\begin{equation}
\int_\Omega\int_{-h/2}^{h/2} \varepsilon_{\alpha\beta}\sigma_{\alpha\beta}\, dx_3 d\Omega + 2\int_\Omega\int_{-h/2}^{h/2} \varepsilon_{\alpha 3}\sigma_{\alpha 3}\, dx_3 d\Omega + \int_\Omega\int_{-h/2}^{h/2} \varepsilon_{33}\sigma_{33}\, dx_3 d\Omega = \delta W_{ext}
\end{equation}

第一项为面内应变项。将总应力 $\sigma_{\alpha\beta}$ 拆分为预应力度与屈曲增量应力：
\begin{equation}
\sigma_{\alpha\beta} = \bar{\sigma}_{\alpha\beta} + \sigma_{\alpha\beta}^{inc}
\end{equation}
其中，预应力度 $\bar{\sigma}_{\alpha\beta}$ 沿厚度均匀分布（与 $x_3$ 无关）；增量应力 $\sigma_{\alpha\beta}^{inc}$ 随厚度弹性变化。将应变展开式代入并完全展开，面内应变分量与应力的乘积共含 8 项：(1) 公式(4.4)中薄膜应变部分与预应力度 $\bar{\sigma}_{\alpha\beta}$ 的乘积；(2) 薄膜应变部分与增量应力 $\sigma_{\alpha\beta}^{inc}$ 的乘积；(3) 广义曲率项 $-x_3\kappa_{\alpha\beta}$ 与预应力度的乘积；(4) 广义曲率项与增量应力的乘积；(5) von K\'arm\'an 项 $\frac{1}{2}w_{,\alpha}w_{,\beta}$ 与预应力度的乘积；(6) von K\'arm\'an 项与增量应力的乘积；(7) 高阶旋转项 $x_3^2\eta_{\alpha\beta}$ 与预应力度的乘积；(8) 高阶旋转项与增量应力的乘积。依据小变形假设和屈曲分岔点的特性，第 1、2 项涉及膜应变，其一阶变分为零；第 6、8 项为增量应力与二阶位移梯度的乘积，属三阶小量舍去；第 3 项因预应力度沿厚度均匀分布与 $x_3$ 无关，$\int x_3 dx_3 = 0$ 恒消去。保留第 4、5、7 项，对厚度积分后定义弯矩合力：
\begin{equation}
M_{\alpha\beta} = \int_{-h/2}^{h/2} x_3 \sigma_{\alpha\beta}^{inc}\, dx_3
\end{equation}
第一项积分的最终结果为：
\begin{equation}
\int_\Omega\int_{-h/2}^{h/2} \varepsilon_{\alpha\beta}\sigma_{\alpha\beta}\, dx_3 d\Omega = \int_\Omega \kappa_{\alpha\beta}M_{\alpha\beta}\, d\Omega + \int_\Omega w_{,\alpha}(h\bar{\sigma}_{\alpha\beta})w_{,\beta}\, d\Omega + \int_\Omega \varphi_{\gamma,\alpha}\left(\frac{h^3}{12}\bar{\sigma}_{\alpha\beta}\right)\varphi_{\gamma,\beta}\, d\Omega
\end{equation}

第二项为横向剪切项。对于纯面内加载的板，预屈曲横向剪应力 $\bar{\sigma}_{\alpha 3} = 0$，故 $\sigma_{\alpha 3} = \sigma_{\alpha 3}^{inc}$。由虎克定律，引入剪切修正系数 $\kappa$ 后：
\begin{equation}
\sigma_{\alpha 3}^{inc} = \kappa G(w_{,\alpha} - \varphi_\alpha)
\end{equation}
代入并对厚度积分（$\int_{-h/2}^{h/2} dx_3 = h$），注意工程剪切应变 $\gamma_\alpha = 2\varepsilon_{\alpha 3}$，第二项的积分结果为：
\begin{equation}
2\int_\Omega\int_{-h/2}^{h/2} \varepsilon_{\alpha 3}\sigma_{\alpha 3}\, dx_3 d\Omega = \int_\Omega \kappa G h\,(w_{,\alpha} - \varphi_\alpha)(w_{,\alpha} - \varphi_\alpha)\, d\Omega
\end{equation}

第三项为厚度方向法向应变项，在平面应力假设 $\sigma_{33}=0$ 下贡献为零。

将三项合并，得到屈曲问题的完整虚功方程：
\begin{equation}
\begin{aligned}
\delta W_{ext} =& \int_\Omega \kappa_{\alpha\beta}M_{\alpha\beta}\, d\Omega + \int_\Omega \kappa G h\,(w_{,\alpha} - \varphi_\alpha)(w_{,\alpha} - \varphi_\alpha)\, d\Omega \\
&+ \int_\Omega h\bar{\sigma}_{\alpha\beta} w_{,\alpha}w_{,\beta}\, d\Omega + \int_\Omega \frac{h^3}{12}\bar{\sigma}_{\alpha\beta} \varphi_{\gamma,\alpha}\varphi_{\gamma,\beta}\, d\Omega
\end{aligned}
\end{equation}
其中前两项为弹性内力虚功（弯曲与剪切），后两项为预应力度对应的几何刚度虚功。当几何刚度项的绝对值等于弹性刚度项时，系统总刚度矩阵奇异，达到临界屈曲状态。

\section{弱形式}

对上述屈曲虚功方程各变量取变分并令其为零，将离散场代入后得到离散形式的广义特征值方程。位移场和转角场的形函数插值为：
\begin{equation}
w(\mathbf{x}) = N_I w_I, \qquad \varphi_\alpha(\mathbf{x}) = N_I \varphi_{\alpha I}
\end{equation}
对空间坐标求偏导，微分算子直接作用在形函数上：$w_{,\alpha} = N_{I,\alpha}w_I$，$\varphi_{\alpha,\beta} = N_{I,\beta}\varphi_\alpha$。虚位移做同阶离散近似。

材料刚度矩阵子块展开如下：面外材料刚度矩阵 $k_b$ 由弯曲刚度（Bending）与横向剪切刚度（Shear）共同组成。对应自由度 $[w,\varphi_1,\varphi_2]^T$，展开后的对称矩阵形式如下：
\begin{equation}
k_b = \begin{bmatrix}
k_{ww} & k_{w\varphi_1} & k_{w\varphi_2} \\
k_{\varphi_1 w} & k_{\varphi_1\varphi_1} & k_{\varphi_1\varphi_2} \\
k_{\varphi_2 w} & k_{\varphi_2\varphi_1} & k_{\varphi_2\varphi_2}
\end{bmatrix}
\end{equation}

各子块的积分表达式为：对于 $k_{ww}$，
\begin{equation}
k_{ww} = \int \kappa Gh(w_{,\alpha}w_{,\alpha})d\Omega
\end{equation}
\begin{equation}
k_{ww} = \int \kappa Gh(N_{I,\alpha}w_I\cdot N_{J,\alpha}w_J)d\Omega
\end{equation}
\begin{equation}
k_{ww} = w_I\cdot\left[\int \kappa Gh(N_{I,\alpha}N_{J,\alpha})d\Omega\right]\cdot w_J
\end{equation}
\begin{equation}
\boxed{K_{IJ}^{ww} = \int \kappa Gh(N_{I,\alpha}N_{J,\alpha})d\Omega}
\end{equation}

对于 $k_{w\varphi_\gamma}$，
\begin{equation}
k_{w\varphi_\gamma} = -\int \kappa Gh(w_{,\alpha}\varphi_\alpha)d\Omega
\end{equation}
\begin{equation}
k_{IJ}^{w\varphi_\gamma} = -\int \kappa Gh(N_{I,\alpha}w_I\cdot N_{\alpha\gamma}N_J\varphi_\gamma)d\Omega
\end{equation}
\begin{equation}
k_{IJ}^{w\varphi_\gamma} = w_I\cdot\left[-\int \kappa Gh(N_{I,\gamma}N_J)d\Omega\right]\cdot\varphi_\gamma
\end{equation}
\begin{equation}
\boxed{k_{IJ}^{w\varphi_\gamma} = -\int \kappa Gh(N_{I,\gamma}N_J)d\Omega}
\end{equation}

对于 $k_{\varphi_\alpha\varphi_\gamma}$，
弯矩合力与曲率的关系为 $M_{\alpha\beta}=D_{\alpha\beta\gamma}\kappa_\gamma$，$\kappa_{\alpha\beta}=\frac{1}{2}(\varphi_{\alpha,\beta}+\varphi_{\beta,\alpha})$。
\begin{equation}
k_{\alpha\gamma}^{\varphi\varphi} = \int(\phi_{\alpha,\beta}D_{\alpha\beta\gamma\eta}\phi_{\gamma,\eta} + \kappa Gh\phi_\alpha\phi_\alpha)d\Omega
\end{equation}
\begin{equation}
k_{I\alpha,J\gamma}^{\varphi\varphi} = \phi_\alpha\cdot\left[\int(N_{I,\beta}D_{\alpha\beta\gamma}N_{J,\eta} + \kappa GhN_IN_J\delta_{\alpha\gamma})d\Omega\right]\cdot\phi_\gamma
\end{equation}
\begin{equation}
\boxed{k_{I\alpha,J\gamma}^{\varphi\varphi} = \int(N_{I,\beta}D_{\alpha\beta\gamma\eta}N_{J,\eta} + \kappa GhN_IN_J\delta_{\alpha\gamma})d\Omega}
\end{equation}
（其中 $\kappa$ 为剪切修正系数，通常薄板或中厚板取 $5/6$）

几何刚度矩阵子块展开如下：展开后的对称矩阵形式为：
\begin{equation}
k_g = \begin{bmatrix}
k_g^{ww} & & \\
& k_g^{\varphi_1\varphi_1} & k_g^{\varphi_1\varphi_2} \\
& k_g^{\varphi_2\varphi_1} & k_g^{\varphi_2\varphi_2}
\end{bmatrix}
\end{equation}

对于 $k_g^{ww}$，
\begin{equation}
k_g^{ww} = \int w_{,\alpha}(h\bar{\sigma}_{\alpha\beta})w_{,\beta}d\Omega
\end{equation}
\begin{equation}
k_g^{ww} = \int(N_{I,\alpha}w_I)\cdot(h\bar{\sigma}_{\alpha\beta})\cdot(N_{J,\beta}w_J)d\Omega
\end{equation}
\begin{equation}
k_g^{ww} = w_I\cdot\left[\int h\bar{\sigma}_{\alpha\beta}N_{I,\alpha}N_{J,\beta}d\Omega\right]\cdot w_J
\end{equation}
\begin{equation}
\boxed{k_{g,IJ}^{ww} = \int h\bar{\sigma}_{\alpha\beta}N_{I,\alpha}N_{J,\beta}d\Omega}
\end{equation}

对于 $k_g^{\varphi_\alpha\varphi_\gamma}$，
\begin{equation}
k_{g,\alpha\gamma}^{\varphi\varphi} = \int \varphi_{,\beta}\left(\frac{h^3}{12}\bar{\sigma}_{\beta}\right)\varphi_{,\gamma}d\Omega
\end{equation}
\begin{equation}
k_{g,I\alpha,J\gamma}^{\varphi\varphi} = \int(N_{J,\beta}\varphi_{\alpha,\alpha})\cdot\left(\frac{h^3}{12}\bar{\sigma}_{\beta}\right)\cdot(N_{I,\varphi_{\gamma,\gamma}})d\Omega
\end{equation}
\begin{equation}
k_{g,I\alpha,J\gamma}^{\varphi\varphi} = \phi_\alpha\cdot\left[\int\frac{h^3}{12}\bar{\sigma}_{\beta}N_{I,\beta}N_{J,\gamma}d\Omega\right]\cdot\phi_\gamma
\end{equation}
\begin{equation}
\boxed{k_{g,I\alpha,J\gamma}^{\varphi\varphi} = \int\frac{h^3}{12}\bar{\sigma}_{\beta}N_{I,\beta}N_{J,\gamma}d\Omega}
\end{equation}

将上述所有子块合并，针对中厚板面外屈曲模式，最终的广义特征值矩阵方程为：
\begin{equation}
\begin{pmatrix}
k_{ww} & k_{w\varphi_1} & k_{w\varphi_2} \\
k_{\varphi_1 w} & k_{\varphi_1\varphi_1} & k_{\varphi_1\varphi_2} \\
k_{\varphi_2 w} & k_{\varphi_2\varphi_1} & k_{\varphi_2\varphi_2}
\end{pmatrix}
+ \lambda
\begin{pmatrix}
k_g^{ww} & & \\
& \frac{h^2}{12}k_g^{ww} & \\
& & \frac{h^2}{12}k_g^{ww}
\end{pmatrix}
\begin{pmatrix} w \\ \varphi_1 \\ \varphi_2 \end{pmatrix}
= \begin{pmatrix} 0 \\ 0 \\ 0 \end{pmatrix}
\end{equation}

% TODO: 待补充——混合离散格式的LBB稳定性分析
% 如需补充，应结合第二章的约束比理论，说明本格式的剪切应力场插值阶数满足inf-sup条件
% 可参考徐洋涛论文第3章的写法：独立成节，含特征值问题推导和约束比数值验证

